\section{Introduction}\label{sec:introduction}
Industrial Control Systems (ICS) are the technological foundation of critical infrastructures such as energy, manufacturing, transportation, and water treatment. The integration of Operational Technology (OT) with Information Technology (IT) and Industrial IoT improved connectivity and efficiency, but it also exposed industrial networks, devices, and supervisory systems to cyber threats. Modern ICS rely on Ethernet and TCP/IP industrial protocols that facilitate interoperability and remote access and introduce new vulnerabilities \cite{Gillen2020ICS,Tama2022IDS}.

Denial-of-Service (DoS) attacks are among the most critical threats, because they disrupt the availability of industrial services by overwhelming devices or networks with malicious traffic. The interruption of the communication between Programmable Logic Controllers (PLCs) \cite{Raman2021ICS}, Human Machine Interfaces (HMIs), and SCADA systems can cause operational instability, safety risks, and economic losses. Unlike conventional IT systems, industrial networks require high availability and real-time responsiveness, which makes them particularly vulnerable to flooding \cite{Benka2025ICS}. PLCs are among the most critical components because they interact directly with the physical process, yet many were designed for reliability rather than security, with limited authentication and encryption \cite{Yu2019IDS}, and protocols such as Modbus/TCP, Profinet, and EtherNet/IP were developed without integrated security \cite{Yu2019IDS}.

Intrusion Prevention Systems (IPS) combine traffic inspection, signature-based detection, anomaly analysis, and automated response to block malicious activity before it reaches industrial operations \cite{Tama2022IDS,Varghese2022DigitalTwinIDS}. However, many studies focus on theoretical analyses, machine-learning models, or intrusion detection without implementing an inline prevention architecture \cite{Makrakis2021ICS,Bhamare2020ICS}, and few measure experimentally what an inline IPS does to PLC availability under DoS attacks.

The objective of this research is to design, implement, and experimentally evaluate an inline Intrusion Prevention System (IPS) that protects the availability of a PLC placed behind a VLAN-segmented virtual industrial network under four flooding attacks. The evaluation is built around a controlled comparison of four conditions (no attack, firewall and VLAN segmentation only, Suricata in alert mode, and Suricata in drop mode) so that the effect of the IPS can be separated from that of the firewall and the segmentation.

The main contributions of this work are:
\begin{itemize}
    \item A reproducible testbed built only with free software (QEMU/KVM, Open vSwitch, nftables, Suricata, OpenPLC) and a scripted run protocol, with the code and the raw measurements publicly available (see Data and Code Availability).
    \item An experimental comparison of the four conditions under SYN, UDP, ICMP, and Modbus/TCP floods, with 10 repetitions per combination (130 runs), 95\,\% confidence intervals, and non-parametric tests.
    \item Metric definitions that separate legitimate from malicious traffic, and a second, independent source (traffic mirrors and packet-count conservation) for every claim about attack traffic removed by the IPS.
    \item Findings that limit the benefit of inline prevention in this testbed: the IPS becomes the bottleneck under volumetric floods, the Modbus/TCP flood crashes the PLC runtime in most runs, and a threshold-based rule produces false positives on the legitimate SCADA client.
\end{itemize}

The paper is organized as follows: Section~\ref{sec:background} gives the background and related work, Section~\ref{sec:method} the method and metrics, Section~\ref{sec:case-study} the architecture and the experimental setup, Section~\ref{sec:results} the results, and Section~\ref{sec:conclusions} the conclusions; the data and code availability closes the paper.
